% ID: SEL-D0056-OB-E04
% ESTADO: cotejada. Segunda pasada visual y borrador compilado revisados.
% Fuente: Selectividad/sel_2011_matematicas_aplicadas/123-EXAMEN 1.pdf
% SHA-256: 69f8b84ed22f2b49074eafb95dbdce8a98114a4a2f02989b92f3f7add9355155
% Página: 2; opción B; ejercicio original 4.

Se sabe que la estatura de las personas de una población es una variable aleatoria que sigue una distribución Normal cuya desviación típica es de 0.04 m. Para estimar la media de esta variable se ha tomado una muestra aleatoria de 60 personas de esa población y se ha encontrado una estatura media de 1.73 m.
\begin{enumerate}[label=\alph*)]
\item \textbf{(1.25 puntos)} Obtenga un intervalo de confianza, con un nivel del 97\%, para la media de la distribución de estaturas.
\item \textbf{(1.25 puntos)} Halle el tamaño mínimo que debe tener una muestra de esta población, para que la amplitud de un intervalo de la media con este nivel de confianza sea inferior a 0.08 m.
\end{enumerate}
